% !TeX program = xelatex
% Contoh pemakaian tema-itb.sty. Salin folder ini, ganti isi, lalu kompilasi:
%   latexmk -xelatex contoh.tex      (atau: xelatex contoh.tex dua kali)
\documentclass[aspectratio=169, 10pt]{beamer}
\usepackage[indonesian]{babel}
\usepackage{tema-itb}

\title{Judul Presentasi Anda}
\subtitle{Subjudul, mis.\ tinjauan pustaka sistematis (Kitchenham), 2021--2026}
\author[Nama Anda]{Nama Anda \textperiodcentered\ NIM \\[1mm]
  {\footnotesize Promotor: Prof.\ Dr.\ Nama Promotor \\
   Ko-promotor: Dr.\ Nama Ko-promotor}}
\institute{Program Doktor Teknik Elektro dan Informatika\\ Sekolah Teknik Elektro dan Informatika, Institut Teknologi Bandung}
\date{Oktober 2026}
% Teks footer bawaan: nama singkat · judul singkat · STEI ITB. Ubah bila perlu:
% \renewcommand{\teksfooter}{Nama Anda \textperiodcentered\ Mata Kuliah \textperiodcentered\ STEI ITB}

\begin{document}
\maketitle

% =====================================================================
\begin{frame}{Agenda}
  \centering
  \begin{tikzpicture}
  \foreach \i/\t/\x/\y in {1/Latar belakang/0/0, 2/Masalah dan pertanyaan penelitian/0/-1.1, 3/Metode/0/-2.2,
                           4/Hasil RQ1/0/-3.3, 5/Hasil RQ2/7.2/0, 6/Hasil RQ3/7.2/-1.1,
                           7/Celah riset dan model/7.2/-2.2, 8/Validitas dan langkah berikutnya/7.2/-3.3} {
    \node[kotakpucat] at (\x,\y) {\t};
    \node[circle, fill=itbBiru, text=white, font=\small\bfseries, minimum size=0.62cm] at (\x-3.2,\y) {\i};
  }
  \end{tikzpicture}
\end{frame}

% =====================================================================
\section{Latar belakang}

\begin{frame}{Slide berisi butir}
  \begin{itemize}
    \item Butir pertama dengan \istilah{istilah kunci} dan \asing{istilah asing}.
    \item Butir kedua dengan \alert{penekanan merah}.
    \begin{itemize}
      \item Sub-butir memakai tanda emas.
    \end{itemize}
    \item Butir ketiga, cukup satu gagasan per butir.
  \end{itemize}
  \sumber{Penulis, Tahun; Penulis lain, Tahun}
  \begin{intinya}
    Satu kalimat yang merangkum pesan utama slide ini.
  \end{intinya}
\end{frame}

\begin{frame}{Slide dua kolom dengan kartu}
  \begin{columns}[T]
    \column{0.5\textwidth}
      \kartu{\linewidth}{Topik}{Isi kartu pertama. Pakai \textbf{tebal} untuk kata yang penting.}
      \kartu[itbEmas]{\linewidth}{Kondisi saat ini}{Kartu dengan kepala emas untuk konteks atau catatan.}
    \column{0.5\textwidth}
      \kartu[itbBiruTua]{\linewidth}{Pertanyaan penelitian}{%
        \chip{1}{Pertanyaan pertama.}\\[2mm]
        \chip{2}{Pertanyaan kedua.}\\[2mm]
        \chip{3}{Pertanyaan ketiga.}}
      \kartu[itbMerah]{\linewidth}{Keterbatasan}{Kartu merah untuk risiko atau keterbatasan.}
  \end{columns}
\end{frame}

\begin{frame}{Temuan utama sekilas}
  \centering
  \begin{tikzpicture}[x=3.35cm, y=2.9cm]
    \node[tile] at (0,0) {\tilenum{22}{104}\\[4pt]\scriptsize makalah primer terpilih};
    \node[tile=itbBiruTua] at (1,0) {\tilenum{22}{12}\\[4pt]\scriptsize label angka kedua};
    \node[tile] at (2,0) {\tilenum{22}{6\,/\,17}\\[4pt]\scriptsize rasio atau perbandingan};
    \node[tile=itbMerah] at (3,0) {\tilenum{22}{0}\\[4pt]\scriptsize temuan negatif};
    \node[tile=itbBiruTua] at (0,-1) {\tilenum{22}{27}\\[4pt]\scriptsize label angka};
    \node[tile] at (1,-1) {\tilenum{22}{3}\\[4pt]\scriptsize label angka};
    \node[tile=itbBiruTua] at (2,-1) {\tilenum{15}{USD\,7,3 jt}\\[4pt]\scriptsize nilai lebar, huruf lebih kecil};
    \node[tile=itbEmas] at (3,-1) {\tilenum{22}{8}\\[4pt]\scriptsize celah riset};
  \end{tikzpicture}
\end{frame}

% =====================================================================
\section{Metode}

\begin{frame}{Diagram alur}
  \centering
  \begin{tikzpicture}
    \node[tahap] (a) at (0,0) {Klarifikasi\\ Riset};
    \node[tahap=itbBiru!55] (b) at (3.6,0) {Studi\\ Deskriptif I};
    \node[tahap=itbBiru!55] (c) at (7.2,0) {Studi\\ Preskriptif};
    \node[tahap=itbBiru!55] (d) at (10.8,0) {Studi\\ Deskriptif II};
    \foreach \p/\q in {a/b,b/c,c/d} \draw[panah] (\p) -- (\q);
    \foreach \x/\s in {0/Semester 1, 3.6/Semester 2, 7.2/Semester 3--4, 10.8/Semester 5--6}
      \node[font=\scriptsize, text=itbAbu] at (\x,-0.95) {\s};
    \draw[itbEmas, line width=2pt, rounded corners=4pt] (-1.75,0.8) rectangle (1.75,-1.25);
    \node[fill=itbEmas, text=black, font=\scriptsize\bfseries, rounded corners=2pt] at (0,1.05) {tahap ini};
  \end{tikzpicture}
  \vspace{4mm}
  \begin{columns}[T]
    \begin{column}{0.32\textwidth}\kartu{\linewidth}{Luaran 1}{Penjelasan luaran pertama.}\end{column}
    \begin{column}{0.32\textwidth}\kartu{\linewidth}{Luaran 2}{Penjelasan luaran kedua.}\end{column}
    \begin{column}{0.32\textwidth}\kartu{\linewidth}{Luaran 3}{Penjelasan luaran ketiga.}\end{column}
  \end{columns}
\end{frame}

\begin{frame}{Tabel}
  \centering\small
  \begin{tabular}{@{}lrrl@{}}
    \toprule
    \textbf{Tahap} & \textbf{Jumlah} & \textbf{Sisa} & \textbf{Catatan} \\
    \midrule
    Pencarian basis data & 1.250 & 1.250 & Scopus, IEEE, ACM \\
    Saring judul/abstrak & $-$980 & 270 & kriteria inklusi \\
    Baca teks lengkap & $-$166 & \textbf{104} & penilaian kualitas \\
    \bottomrule
  \end{tabular}
\end{frame}

\begin{frame}[standout]
  Terima kasih\\[4mm]
  {\normalsize Nama Anda \textperiodcentered\ NIM\\[1mm]
  \footnotesize nim@mahasiswa.itb.ac.id \textperiodcentered\ situs-anda.id}\\[5mm]
  \begin{tikzpicture}\fill[white] (0,0) circle (0.8cm);\node at (0,0) {\includegraphics[height=1.4cm]{\logoitb}};\end{tikzpicture}
\end{frame}

\end{document}
